\chapter{Python、NumPy 与 Pandas 速查}
本附录提供本书实验的最小预备知识。Python 在科学计算中承担数据读取、数组运算、训练与可视化。学生应能阅读变量赋值、函数调用、条件和循环，不需要先掌握完整的软件工程体系。
\section{列表与数组}
Python 列表可以保存多种对象。NumPy是Python数值计算软件包，数组是它用来存放规则数值集合的结构；程序中的np只是import时约定的简写。数组的形状表示各个轴的长度：一组100条、每条5个特征的数据通常形状是 $(100,5)$。数组 \texttt{X[:, 0]} 表示第一列；\texttt{X[0, :]} 表示第一行。监督学习目标常用一维数组 \texttt{y} 表示，长度应与样本行数一致。
\begin{lstlisting}[caption={数组的形状、切片与基本统计}]
import numpy as np
X = np.array([[1.0, 2.0], [2.0, 4.0], [3.0, 8.0]])
y = np.array([0, 0, 1])
print(X.shape, y.shape)
print(X[:, 0], X.mean(axis=0))
\end{lstlisting}
\section{DataFrame 的核心操作}
Pandas是处理带行列标签表格的软件包，通常简写为\texttt{pd}。它的\texttt{DataFrame}是二维表，\texttt{Series}是一列或一组带索引的值。CSV是用分隔符组织字段的文本格式，\texttt{read\_csv}把文件读成表；它不会自动判断哪些字段是科学输入或目标。

\texttt{DataFrame} 在保留列名、索引与类型方面非常方便。\texttt{df.loc} 按行列标签选择，\texttt{df.iloc} 按位置选择；\texttt{groupby} 用于分组统计；\texttt{merge} 可以连接表格，但必须检查连接键的唯一性与行数变化。科学对象的 ID 不应因为参与连接而被转换为有损浮点数。
\begin{lstlisting}[caption={CSV 读取与字段检查}]
import pandas as pd
df = pd.read_csv("observations.csv")
print(df.head(), df.dtypes)
print(df.isna().sum())
# Only select fields that are valid features at prediction time.
X = df[["feature_1", "feature_2"]]
\end{lstlisting}
上面的字段名仅为模板，占位字段必须换成真实数据字典中的列。项目保存 CSV 时建议显式说明是否保留行索引；不要将 Pandas 自动索引与科学对象身份混为一谈。
\section{绘图与随机数}
Matplotlib 常用于绘制散点、直方图、损失曲线、混淆矩阵与残差。每幅科研图至少应有坐标名称、单位、样本范围和必要的图例。随机模拟使用 NumPy 的 \texttt{default\_rng(seed)}，训练与划分使用 sklearn 的 \texttt{random\_state}。种子可复现特定运行，但不代表结果在任何随机状态下都稳定。
\begin{lstlisting}[caption={最小绘图示例}]
import matplotlib.pyplot as plt
fig, ax = plt.subplots()
ax.scatter([1, 2, 3], [2, 2.5, 4])
ax.set(xlabel="Input (unit)", ylabel="Target (unit)")
fig.tight_layout()
fig.savefig("scatter.pdf")
\end{lstlisting}
\section{Notebook 与脚本}
Jupyter Notebook是一种把文字、代码和运行输出组织在多个单元中的交互文档。\textbf{内核}是实际执行代码并保存变量的进程。单元在页面上的先后不保证执行先后；因此后写的代码可能依赖一个已经在内存中、却没有出现在当前前文的变量。重启内核后从头运行，正是为了检查这类隐藏依赖。

Notebook 适合交互探索，但发布前必须重启内核顺序执行，避免隐藏状态。数据下载、清洗和核心图表生成最好也保留普通 Python 脚本，便于批量核验。使用可复现路径和环境文件，避免在代码中写死作者个人电脑的绝对路径。


\section{从一行记录到一批记录的形状}
单个对象的13项特征在数学上可以写成一个13维向量，但多数sklearn预测接口期望二维矩阵：行是样本，列是特征。因此一条新记录应组织为形状 $(1,13)$，而不是单纯长度13的数组。Pandas的 \texttt{df.iloc[[0]]}保留二维表，\texttt{df.iloc[0]}通常返回一维Series。理解这个区别，可以避免为了消除报错而随意转置。

\begin{lstlisting}[caption={保持特征顺序和二维样本结构}]
columns = list(X_train.columns)
new_rows = incoming.loc[:, columns].copy()
assert new_rows.shape[1] == X_train.shape[1]
assert not new_rows.columns.duplicated().any()
prediction = fitted_pipeline.predict(new_rows)
\end{lstlisting}

列名一致还不保证单位一致。若训练记录使用厘米，新记录使用毫米，需要先按字段规范转换，不能因为列名相同就认为输入可用。对于字符串标识，保持原样用于连接和结果追踪，不为了让所有列都是浮点而把标识一起转换。

\section{检查连接关系与缺失来源}
\texttt{merge}是按键连接两张表，不是简单按行位置拼接。若主表一个对象一行、辅助表同对象多行，连接会扩展记录数；这是合理的一对多关系还是不应有的重复，要由样本单位决定。
\begin{lstlisting}[caption={给一对一连接设置显式约束}]
assert main_table['source_id'].is_unique
assert auxiliary['source_id'].is_unique
joined = main_table.merge(auxiliary, on='source_id',
                          how='left', validate='one_to_one',
                          indicator=True)
print(joined['_merge'].value_counts())
\end{lstlisting}

没有匹配上的值、原始测量本来缺失和读文件时把字符串解析失败，是三种不同来源。检查缺失时保留来源记录，避免把所有空值统一解释为“随机漏测”。写CSV时指定是否保留索引，读入后核对列名、类型和记录数；不要让自动生成的索引列意外进入模型。

\section{Notebook的执行状态}
Notebook允许任意顺序执行单元，这对探索方便，却可能让后面的结果依赖未保存的旧变量。提交前重启内核并从上到下运行，确认每个变量在使用前定义，每个输入来自说明过的文件或代码。图表输出应由当前源码生成，不保留已经不对应当前代码的旧结果。

临时探索可以另存草稿；正式实验尽量把数据准备、配置、训练、评价和保存按顺序组织。错误消息先检查数据形状、类型、路径和接口参数，再考虑模型是否有统计问题。异常分数高与程序抛出异常并不是同一概念。

